% =====================================================================
%  2bat-ud6-14.tex  —  SESSIÓ 14: Repàs global i síntesi (estacions)
%  (sense preàmbul: s'inclou des de main.tex)
% =====================================================================
\horatitol{Sessió 14 — Repàs global i síntesi}%
{Repàs actiu per quatre estacions, una per bloc de la unitat, i mapa final que connecta mostreig, intervals, significació i lectura crítica.}

\subsection{Idea clau}

Repassem \textbf{tota la unitat} de manera activa, rotant per \textbf{quatre estacions},
una per cada bloc treballat. Al final construirem un \textbf{mapa de la unitat} que connecta
les quatre idees —\textbf{mostreig}, \textbf{intervals de confiança}, \textbf{significació}
i \textbf{lectura crítica}— i, com que és la darrera unitat, hi afegirem una mirada enrere
sobre tot el recorregut de l'etapa.

\begin{idea}
\textbf{Les quatre estacions}
\begin{enumerate}[leftmargin=1.6em, itemsep=2pt]
  \item \textbf{Mostreig i representativitat:} valorar si una mostra és representativa i
        identificar-ne els biaixos.
  \item \textbf{Intervals de confiança:} construir i interpretar l'interval d'un resultat
        amb marge.
  \item \textbf{Significació:} valorar si una diferència entre dos grups és real o casual.
  \item \textbf{Lectura crítica:} analitzar un titular o fragment d'informe i detectar-ne
        els paranys.
\end{enumerate}
\end{idea}

\subsection{Les estacions}

\begin{exemple}[Estació 2 — intervals de confiança]
Resultat: «$44\%\pm 3\%$». Interval: $[41\%,\,47\%]$. Lectura: és molt probable que el valor
real de tota la població estigui entre el $41\%$ i el $47\%$; no exactament el $44\%$.
\end{exemple}

\begin{exemple}[Estació 3 — significació]
Dos grups: $50\%\pm 3\%$ ($[47\%,53\%]$) i $52\%\pm 3\%$ ($[49\%,55\%]$). Els intervals
\textbf{se superposen}: la diferència de $2$ punts \textbf{no} és significativa (pot ser
atzar).
\end{exemple}

\subsection{Exercicis resolts}

\begin{exresolt}[model de l'Estació 1: mostreig]
Una enquesta sobre l'ús d'un servei juvenil es fa \textbf{només per Instagram}. És
representativa? Quin biaix té?
\solucio
\textbf{No} és representativa: arriba només als joves que segueixen el servei a Instagram
(perfil ja implicat i connectat). És un \textbf{biaix de cobertura i de selecció}: en queden
fora els joves que no usen aquesta xarxa o no segueixen el compte. Sobreestimaria l'interès
pel servei.
\resposta{No; biaix de cobertura/selecció, sobreestima l'interès.}
\end{exresolt}

\begin{exresolt}[model de l'Estació 4: lectura crítica]
Titular: «\emph{El servei triomfa: el $90\%$ el recomana}», a partir d'una enquesta feta
només als usuaris actuals. Quin parany hi ha?
\solucio
\textbf{Biaix de selecció}: s'enquesta només qui continua sent usuari (qui n'estava content),
i no qui va marxar. El $90\%$ està inflat i no es pot generalitzar. A més, «triomfa» presenta
com a certesa una dada d'una mostra esbiaixada.
\resposta{Biaix de selecció (només usuaris actuals); el $90\%$ està inflat.}
\end{exresolt}

\subsection{Exercicis per fer}

\noindent\textit{Una tasca per estació. Rota i resol-les totes.}

\begin{exercicis}
\begin{exlist}
  \item \textbf{(Estació 1)} Una enquesta sobre necessitats de la gent gran es fa per
        internet. És representativa? Qui en queda fora?
  \item \textbf{(Estació 2)} Construeix l'interval de «$36\%\pm 4\%$» i escriu què afirma
        correctament l'estudi.
  \item \textbf{(Estació 3)} Dos grups: $60\%\pm 3\%$ i $70\%\pm 3\%$. Se superposen els
        intervals? La diferència és significativa?
  \item \textbf{(Estació 4)} Titular: «\emph{Els que fan esport són més feliços}», a partir
        d'una correlació. Quin parany hi ha?
  \item \textbf{(Mapa de la unitat)} Completa el mapa connectant cada bloc amb la seva idea:
        \begin{center}
        \renewcommand{\arraystretch}{1.4}
        \begin{tabular}{p{4.0cm} p{8.4cm}}
        \toprule
        \textbf{Bloc} & \textbf{Idea clau} \\
        \midrule
        Mostreig i representativitat & \rule{7.6cm}{0.4pt} \\
        Intervals de confiança & \rule{7.6cm}{0.4pt} \\
        Significació & \rule{7.6cm}{0.4pt} \\
        Lectura crítica & \rule{7.6cm}{0.4pt} \\
        \bottomrule
        \end{tabular}
        \end{center}
  \item \textbf{(Mirada enrere)} En una línia, resumeix el recorregut de l'etapa: de
        \underline{\hspace{2cm}} a \underline{\hspace{2cm}} a \underline{\hspace{2cm}} a
        \underline{\hspace{2cm}}.
\end{exlist}
\end{exercicis}

% ---------------------------------------------------------------------
\fullsolucions{Sessió 14}
\begin{solucions}
\begin{sollist}
  \item No: en queden fora les persones grans sense accés o competència digital (sovint les
        més vulnerables). Biaix de cobertura; el resultat descriuria només la gent gran
        connectada.
  \item Interval $[32\%,\,40\%]$. És molt probable que el valor real de la població estigui
        entre el $32\%$ i el $40\%$, no exactament el $36\%$.
  \item Intervals $[57\%,63\%]$ i $[67\%,73\%]$: \textbf{no} se superposen. La diferència de
        $10$ punts és \textbf{significativa}.
  \item Correlació presentada com a causa: fer esport i ser feliç poden anar junts sense que
        l'un causi l'altre (potser qui està millor d'ànim fa més esport, o hi ha una tercera
        causa). El titular sobreinterpreta.
  \item Resposta oberta. Idees: \textbf{mostreig} = de la mostra a la població, cal que
        sigui representativa; \textbf{intervals} = expressar la incertesa amb un marge;
        \textbf{significació} = valorar si una diferència és real o casual; \textbf{lectura
        crítica} = detectar limitacions i paranys en estudis i titulars.
  \item Resposta oberta. Per exemple: de \textbf{descriure} a \textbf{mesurar la incertesa}
        a \textbf{modelar} a \textbf{inferir i criticar}.
\end{sollist}
\end{solucions}
